% Generated draft: requires mathematical and source review.
\documentclass[11pt]{article}
\usepackage[margin=25mm]{geometry}
\usepackage{fontspec}
\setmainfont{Latin Modern Roman}
\usepackage{amsmath,amssymb,mathtools,unicode-math}
\setmathfont{Latin Modern Math}
\usepackage{enumitem,needspace,xurl,hyperref}
\setlength{\parindent}{0pt}
\setlength{\parskip}{5pt}
\setlength{\emergencystretch}{3em}
\newcommand{\N}{\mathbb{N}}
\newcommand{\Z}{\mathbb{Z}}
\newcommand{\Q}{\mathbb{Q}}
\newcommand{\R}{\mathbb{R}}
\newcommand{\C}{\mathbb{C}}
\newcommand{\F}{\mathbb{F}}
\newcommand{\A}{\mathbb{A}}
\newcommand{\PP}{\mathbb{P}}
\newcommand{\E}{\mathbb{E}}
\newcommand{\eps}{\varepsilon}
\newcommand{\dd}{\,\mathrm{d}}
\begin{document}
\section*{UnsolvedMath ID: 800925}
\textbf{Source problem code:} catalogue entry
\subsection{Definitions and mathematical statement}
Let \(M=\R\times S^1\), where \(S^1=\R/\Z\). Under the usual smooth, strongly convex convention, a Finsler metric is a function \(F:TM\to[0,\infty)\) that is continuous everywhere, smooth away from the zero section, positive on nonzero tangent vectors, positively homogeneous,
\[
F(x,\lambda v)=\lambda F(x,v)\qquad(\lambda>0),
\]
and such that the fibrewise Hessian of \(\tfrac12F^2\) is positive definite at every nonzero tangent vector. It is \emph{irreversible} if
\[
\exists x\in M\ \exists v\in T_xM\setminus\{0\}:\quad F(x,v)\ne F(x,-v).
\]
A closed geodesic means a nonconstant periodic \(F\)-geodesic: a geodesic \(\gamma:\R\to M\) for which there is \(T>0\) such that \(\gamma(t+T)=\gamma(t)\) for every \(t\in\R\).

The supplied question, attributed in context to Bangert, asks whether there exists an irreversible Finsler metric \(F\) on \(M\) satisfying
\[
\operatorname{Vol}_{F}(M)<\infty
\qquad\text{and}\qquad
\text{there is no nonconstant closed }F\text{-geodesic}.
\]
The source does not specify which Finsler volume is intended; thus \(\operatorname{Vol}_{F}\) here records an unresolved choice of volume construction. It also does not specify regularity or convexity conventions beyond the term Finsler metric; the definition above states the conventional interpretation used in this draft. No forward or backward completeness hypothesis is stated, and neither is imposed here.

The supplied background expects reversible Finsler metrics to behave similarly to Riemannian metrics for questions of this kind. It cites Katok's example as showing that compact irreversible Finsler manifolds can have only finitely many closed geodesics, and presents as a possibility that every noncompact finite-volume Finsler manifold might nevertheless have infinitely many closed geodesics. These are source-context assertions, not additional hypotheses or independently verified conclusions. The source describes the displayed existence question as completely open; this is a supplied status description.
\subsection{Short English statement}
Does the cylinder \(\R\times S^1\) admit an irreversible Finsler metric of finite total volume with no nonconstant closed geodesics? The intended Finsler volume and precise metric regularity conventions are unspecified, and the source states no completeness assumption. The supplied record labels the question open and associates it with Bangert.
\subsection{Sources}
\begin{itemize}
\item Local UnsolvedMath registry, permanent ID \(800925\), source problem code \texttt{catalogue entry}, titled \emph{Open Problems and Questions about Geodesics --- Question 5 (Bangert)}. The record lists Burns and Matveev as proposers, gives the year 2019, and labels the problem open. These are supplied registry data.
\item Burns and Matveev, \emph{Open Problems and Questions about Geodesics}, listed in the supplied source record as 2019, Question 5. Supplied source URL: \url{https://arxiv.org/abs/1308.5417}. The record reports access on 2026-10-16 and extraction from \texttt{arxiv\_latex}; this access and extraction information is supplied provenance, not an independent source verification.
\item The supplied background gives the attribution/context as Bangert and mentions Katok's example. It supplies no separate bibliographic reference for either attribution or example.
\end{itemize}
\end{document}
